\section{Machine Learning}
\label{sec:ml}
The mass fits separate signal from background through the shape of the mass spectrum. As a complementary, data-driven approach, classifiers were trained to separate them from the other variables of the candidates, and used to test whether a machine-learning selection changes $\Araw$ and to partition the Dalitz plane without fixing the binning in advance.

\subsection{Weak Labels and Classifiers}
\label{sec:ml_labels}
\label{sec:ml_models}
No candidate carries a true label, so the classifiers were trained on weak labels derived from the mass spectrum, following the classification-without-labels approach~\cite{cwola}: a classifier that separates two mixtures with different signal fractions, here the window $|m-\mu|<2\sigma$ ($15{,}130$ signal-like candidates) and the sidebands $[5150,\mu-4\sigma]\cup[\mu+4\sigma,5600]\,\MeV$ ($1940$ background-like), is, by the Neyman-Pearson lemma, also the optimal signal/background classifier, provided that the features carry no information on the mass\appref{sec:ml_labels}. The $20$ features are the kinematic and particle-identification variables of the three tracks and two vertex variables of the $B$; the mass is not an input. Figure~\ref{fig:ml_labels} shows the regions on the mass spectrum. The classes are imbalanced ($0.128$) and are balanced with class weights, and the sample is split $60/20/20$ into training, validation and test sets. Six classifiers were compared: logistic regression with $L2$ and $L1$ penalty, a decision tree, a Random Forest, gradient boosting and a feed-forward neural network (FFNN). The tree ensembles and the network reach a test AUC of $0.857$--$0.864$, $0.09$ above the linear models and $0.04$ above the single tree, and are indistinguishable from each other within the uncertainty of the test set ($0.012$); gradient boosting, the best in cross-validation, is the reference\appref{sec:ml_models}.

%%FIG:fig:ml_labels%%

\subsection{Regularization, Bias-Variance and Uncertainty}
\label{sec:ml_regcal}
\label{sec:ml_uncertainty}
The regularization strength of the logistic regression, scanned over $30$ values of $C=1/\alpha$ with cross-validation, gives the same optimal AUC for the $L1$ and $L2$ penalties ($0.7686$) and no overfitting: the limitation of the linear model is its decision boundary, not the variance. For gradient boosting, whose complexity grows with the number of stages, the test AUC rises up to $\sim150$ stages and then stays at $0.865$ while the training AUC keeps rising (the gap grows from $0.018$ at $20$ stages to $0.084$ at $300$): the additional stages overfit without gain\appref{sec:ml_regcal}. Class weights bias the probabilities, which converge to $\tilde p=w_sp/[w_sp+w_b(1-p)]$ instead of $p$; an isotonic recalibration on held-out folds corrects it (Brier score $0.122\to0.062$) without changing the ranking. The uncertainty of a prediction was estimated from the spread of the $300$ trees of the Random Forest, of $300$ forward passes with dropout kept active (MC dropout) and of a deep ensemble of $10$ networks: the mean per-candidate values are $0.178$, $0.095$ and $0.050$; they measure different sources of variability, and only the Random Forest spread decreases visibly with the size of the training set ($0.220$, $0.204$, $0.179$ with a quarter, half and all of it), the epistemic part, while the part at scores around $0.5$ is the aleatoric one\appref{sec:ml_uncertainty}.

\subsection{Effect on the Measurement and Dalitz Partition}
\label{sec:ml_stability}
\label{sec:ml_kmeans}
A cut at $0.5$ on the score keeps $92\%$ of the fitted signal and $36\%$ of the background ($15{,}709$ candidates). The raw asymmetry is stable (Fig.~\ref{fig:ml_three}): $-0.0451\pm0.0087$ (Bayesian) and $-0.0447\pm0.0086$ (Frequentist), against $-0.0443$ and $-0.0441$ for the full window, all shifts within the spread expected for a subsample and $5.2\sigma$ from zero in both samples. The precision does not improve, because the nominal sample has little background to remove ($\approx4\%$ in the signal window). The mass spectrum, instead, changes ($\chi^2$ from $164/88$ to $253/88$ in the Bayesian fit): the classifier does not sculpt the mass, since without the nine momentum components the degradation is the same ($\chi^2=533/388$ in the Frequentist fit), but the selection changes the composition of the background, keeping partially reconstructed decays that the model does not describe\appref{sec:ml_stability}. Finally, $K$-means with $K=8$ on $(m^2_{\mathrm{lo}},m^2_{\mathrm{hi}})$, an unsupervised partition learned from the density of the candidates and stable between initializations (adjusted Rand index $>0.998$), finds two clusters beyond $3\sigma$ ($\chi^2=46.4$ for $7$ degrees of freedom) with the pattern of the fixed binning: negative at low $m^2_{\mathrm{lo}}$ ($-0.127\pm0.017$) and positive at high $m^2_{\mathrm{hi}}$ (Fig.~\ref{fig:ml_kmeans})\appref{sec:ml_kmeans}.

%%FIG:fig:ml_kmeans%%

%%FIG:fig:ml_three%%
